\documentclass[10pt,twocolumn,letterpaper]{article}

% Standalone one-page evidence-table workspace. It deliberately retains the
% official ACCV rebuttal geometry and typography from rebuttal copy.tex.
\usepackage[rebuttal]{cvpr}
\usepackage{accvabbrv}
\usepackage{graphicx}
\usepackage{booktabs}
\usepackage{array}
\definecolor{accvblue}{rgb}{0.12,0.49,0.85}
\usepackage[pagebackref,breaklinks,colorlinks,citecolor=accvblue]{hyperref}

\def\paperID{1107}
\def\confName{ACCV}
\def\confYear{2026}
\setcounter{table}{5} % Continue after Tables 1--5 in the submission.

\begin{document}
\title{Exploring RGB and Video V-JEPA 2.1 Features for Aerial Depth Estimation and Semantic Scene Completion}
\maketitle
\thispagestyle{empty}
\appendix

\noindent\textbf{We thank all reviewers for their carefully reviewed constructive feedback.}
VJ/D2/D3/DA2/MVSF abbreviate V-JEPA~2.1/DINOv2/DINOv3/Depth Anything~2/MVSFormer++;
response table numbering starts at 6, continuing submission.
\begin{table}[t]
  \centering
  \caption{\textbf{Additional depth controls extending main-paper Table~3.}}
  \label{tab:depth_evaluation}
  \footnotesize
  \setlength{\tabcolsep}{2.5pt}
  \begin{tabular}{@{}lcccc@{}}
    \toprule
    Method & $\delta_1\uparrow$ & AbsRel$\downarrow$ & RMSE$\downarrow$ & SILog$\downarrow$ \\
    \midrule
    VJ DPT          & 0.747 & 0.162 & 6.730 & 0.200 \\
    VJ video Linear Probe              & 0.701 & 0.191 & 7.055 & 0.207 \\
    D3 Linear Probe                    & 0.782 & 0.156 & 6.097 & 0.186 \\
    D3 FiLM-DPT              & \textbf{0.848} & \textbf{0.121} & \textbf{4.358} & 0.161 \\
    VJ DTA-MVS     & 0.654 & 0.231 & 8.197 & 24.717 \\
    VJ DTA-FiLM-MVS          & 0.768 & 0.175 & 6.406 & 0.211 \\
    MVSF++/D2-B             & 0.748 & 0.231 & 7.657 & 0.195 \\
    MVSF++/VJ-G             & 0.761 & 0.225 & 7.517 & 0.193 \\
    MVSF++/D3-H+            & 0.828 & 0.136 & 4.611 & \textbf{0.130} \\
    \bottomrule
  \end{tabular}
\end{table}

\paragraph{Responses common to LD5V, CWVJ, and y7qY:}
Table~\ref{tab:depth_evaluation} adds the corrected DPT result, spatially aligned video methods, DINOV3 comparision, simple-MVS, and supportive sophisticated MVSF++ variants. For space, we omit MAE, $\delta_2$, $\delta_3$, and published baselines already reported in OccuFly Table~5, retaining DA2-OccuFly as the key comparator. We withdraw the VJ FiLM-DPT superiority claim: Although DAv2-OccuFly achieves the strongest reported SILog and threshold accuracies, foundation models with general pretraining objectives remain competitive and outperform it on absolute-error metrics. Any remaining gap may reflect DAv2’s depth-specific pretraining and end-to-end OccuFly adaptation, whereas frozen VJ/D backbones are adapted only through their depth heads.

\begin{table}[t]
  \caption{\textbf{Altitude interpolation diagnostics.} }
  \label{tab:altitude-interpolation-control}
  \centering
  \footnotesize
  \setlength{\tabcolsep}{2.2pt}
  \resizebox{\columnwidth}{!}{%
  \begin{tabular}{@{}lccccc@{}}
    \toprule
    Method & $\delta_1\uparrow$ & AbsRel$\downarrow$ & RMSE$\downarrow$ &
      MAE$\downarrow$ & SILog$\downarrow$ \\
    \midrule
    \multicolumn{6}{@{}l}{Tr: 01--08@30/50\,m; Val: 01--04@40\,m; Te: 05--08@40\,m} \\
    DPT
      & 0.789 & 0.151 & 6.728 & 5.775 & 0.180 \\
    FiLM-DPT
      & \textbf{0.965} & \textbf{0.102} & \textbf{4.707} &
        \textbf{3.957} & \textbf{0.120} \\
    \bottomrule
  \end{tabular}
  }
\end{table}

Table~\ref{tab:altitude-interpolation-control} shall be Table~4. This non-standard split isolates interpolation to an unseen altitude while retaining familiar scenes; Suggested leakage is real and hence a new ablation with unseen scene and altitude will be added. We will replace Figures~1 and~5 with a clear end-to-end SSC data-flow diagram.


\begin{table}[t]
  \caption{\textbf{FiLM conditioning ablation.} Test mean$\pm$SD over three
  seeds; best mean in \textbf{bold}.}
  \label{tab:film-aggregate}
  \centering
  \footnotesize
  \setlength{\tabcolsep}{2.6pt}
  \resizebox{\columnwidth}{!}{%
  \begin{tabular}{@{}lcccc@{}}
    \toprule
    Conditioning & $\delta_1\uparrow$ & AbsRel$\downarrow$ & RMSE$\downarrow$ & SILog$\downarrow$ \\
    \midrule
    DPT
      & $.728{\pm}.034$ & $.174{\pm}.006$ & $6.709{\pm}.385$ & $.202{\pm}.010$ \\
    Altitude FiLM
      & $.798{\pm}.020$ & $.145{\pm}.005$ & $5.236{\pm}.322$ & $.180{\pm}.006$ \\
    Intrinsics FiLM
      & $.711{\pm}.015$ & $.182{\pm}.004$ & $7.115{\pm}.243$ & $.216{\pm}.007$ \\
    Alt.+intr. FiLM
      & $\mathbf{.810{\pm}.006}$ & $\mathbf{.137{\pm}.005}$ & $\mathbf{4.907{\pm}.140}$ & $\mathbf{.171{\pm}.005}$ \\
    \bottomrule
  \end{tabular}
  }
\end{table}
\paragraph{LD5V:}
Spatial alignment: Frozen RAFT estimates target-to-source flow and warps neighboring 2D features into target-frame coordinates. For a stronger MVS comparison, MVSF++ uses cross-view fusion, attention-based encoding, and cost-volume regularization. Extending main-paper Table~5, its D3 and VJ variants integrated into FSSC yield SC IoU/SSC mIoU of 43.37/4.47\% and 35.02/3.28\%, respectively. Replacing estimated depth with GT metric depth only during geometric lifting yields 49.40/7.25\% with VJ context and 42.01/8.40\% with D3 context. Per-class results are provided in Supplement~B; MVSF++-FSSC variants will be added.

\paragraph{CWVJ:}
Suggested D3 comparision also motivates, if permitted, a revised title:
``Exploring Foundation-Model Spatial and Temporal Features for Aerial Depth Estimation and Semantic Scene Completion.''
Table \ref{tab:film-aggregate} for FiLM contribution ablations given the limited rebuttal period, and can extend this analysis to all published results.
Our study focuses on frozen-feature transfer for MDE and SSC rather than onboard deployment; we therefore defer SWaP profiling to compact UAV-suitable student models and will add the suggested related work. Existing aerial datasets with metric depth can extend MDE evaluation, but those without voxel-level annotations cannot support SSC evaluation, for which OccuFly provides the required ground-truth voxel grids.


\paragraph{y7qY:}
Low SSC mIoU persists with GT depth (7.25/8.40), indicating bottlenecks beyond depth; learning is further challenged by OccuFly’s 96.109\% empty versus 3.891\% occupied valid voxels. The video result shows no benefit from our current adapter—not that temporal cues are ineffective—motivating motion/pose-aligned, occlusion-aware fusion. MVSFormer++ provides stronger cross-view geometry than our simple MVS module. The altitude-interpolation ablation uses scenes 1–8 only.

\end{document}
